% R32 section 2: literature review (input by R32_floquet_spectra.tex)
\section{The tools, and what they mean for a Floquet operator}
\label{sec:lit}

Everything in this report is a property of one operator, the one-period
evolution $U=L_BL_A$ of the brickwork circuit, and not of a Hamiltonian. This
section collects what the literature says the three diagnostics mean in that
setting, and where the Hamiltonian intuition has to be dropped.

\subsection{Quasienergies, and why there is no ``middle of the spectrum''}

For a time-periodic Hamiltonian $H(t+T)=H(t)$ or a periodically applied circuit,
the stroboscopic dynamics is generated by the Floquet operator
$U=\mathcal{T}\exp[-i\int_0^T H(t)\,dt]$, here simply the product of the two
gate layers. Its eigenvalues lie on the unit circle,
$U\ket{\phi_n}=e^{-i\theta_n}\ket{\phi_n}$, and the quasienergies $\theta_n$ are
defined only modulo $2\pi$. There is no ordering, no ground state and no
temperature axis.

For a generic interacting driven system the long-time state is the
\emph{infinite-temperature} state. D'Alessio and Rigol showed by exact
diagonalisation and finite-size scaling that for any non-zero driving period the
Floquet operator of a generic lattice system acquires circular-ensemble
properties as the system grows. The Floquet Hamiltonian is then non-local and
the system heats to infinite temperature~\cite{dalessio2014}. Lazarides, Das and
Moessner reached the same conclusion from the structure of Floquet eigenstates,
which mix many eigenstates of the undriven Hamiltonian. They also showed that
integrable driven systems are instead described by a periodic generalised Gibbs
ensemble~\cite{lazarides2014}. Ponte, Chandran, Papi\'c and Abanin found that
ergodic systems delocalise in energy space and heat up under both local and
global driving~\cite{ponte2015annphys}. Disorder can prevent this: the Floquet
eigenstates of a many-body localised system obey an area law, and a delocalised
phase is separated from the MBL phase by a transition seen in the eigenstate
entanglement~\cite{ponte2015prl}.

\emph{Consequence for the figures.} The Floquet version of the eigenstate
thermalisation hypothesis
(ETH)~\cite{deutsch1991,srednicki1994,rigol2008,dalessio2016} says that every
eigenstate of a chaotic $U$ looks locally like the infinite-temperature state.
Plotted against $\theta$, the eigenstate entanglement of a thermalising circuit
is therefore \emph{flat}, a band at the random-state value at every
quasienergy. There is no Hamiltonian ``dome'' that peaks in the middle of the
spectrum. A non-thermal (scar-like) eigenstate shows up as an outlier
\emph{below} the band, at whatever quasienergy it happens to sit.

\subsection{Spectral statistics: circular ensembles and the gap ratio}

Dyson's circular ensembles are the random-matrix models for unitary
operators~\cite{dyson1962}: orthogonal (COE, $\beta=1$), unitary (CUE,
$\beta=2$) and symplectic (CSE, $\beta=4$). The Bohigas--Giannoni--Schmit
conjecture~\cite{bgs1984} assigns random-matrix level statistics to quantum
chaotic systems. The Berry--Tabor conjecture assigns Poisson statistics to
integrable ones, with uncorrelated, clustering levels~\cite{berrytabor1977}.
Which circular ensemble applies is fixed by the antiunitary symmetries \emph{of
$U$ itself}. This need not match the class of the instantaneous Hamiltonian:
Regnault and Nandkishore showed that a Floquet system can be in the unitary
class while its instantaneous Hamiltonian is always orthogonal, and vice versa.
They also showed that a drive can thermalise even when the instantaneous
Hamiltonian is integrable~\cite{regnault2016}.

The gap ratio $r_n=\min(s_n,s_{n+1})/\max(s_n,s_{n+1})$ of consecutive
spacings, introduced by Oganesyan and Huse~\cite{oganesyan2007}, is independent
of the local density of states and so needs no unfolding. Atas \emph{et al.}\
derived Wigner-like surmises for its distribution
\begin{equation}
  P_\beta(r)=\frac{2}{Z_\beta}\,\frac{(r+r^2)^\beta}{(1+r+r^2)^{1+3\beta/2}},
  \qquad r\in[0,1],
  \label{eq:surmise}
\end{equation}
with $Z_1=8/27$ and $Z_2=4\pi/(81\sqrt3)$. Their large-matrix values,
$\langle r\rangle=0.5307$ (GOE) and $0.5996$ (GUE), are the reference values
used here. Poisson statistics give $P_0(r)=2/(1+r)^2$ and
$\langle r\rangle=2\ln2-1=0.3863$~\cite{atas2013}. For a chaotic Floquet
operator the quasienergy density is uniform, so the ratio is doubly convenient.

Two caveats from the literature are central to this report:
\begin{itemize}
\item \textbf{Symmetry sectors must be resolved.} Superposing independent
  spectra from different symmetry sectors removes level repulsion between them
  and pulls $\langle r\rangle$ towards the Poisson value, even if every sector
  is chaotic. Giraud, Mac\'e, Vernier and Alet derived the gap-ratio surmises for
  such superpositions and used them to detect hidden symmetries, including in
  periodically driven spin chains~\cite{giraud2022}. ETH, too, is a statement
  per symmetry sector. Mondaini \emph{et al.}\ made exactly this point in their
  comment~\cite{mondaini2018} on a claimed ETH counterexample~\cite{shiraishi2017}.
\item \textbf{In fragmented systems the Krylov sectors play the role of
  symmetry sectors.} Hilbert-space fragmentation (HSF) splits the space into
  exponentially many invariant subspaces that are not labelled by conventional
  symmetries~\cite{sala2020,khemani2020}. In the commutant language of Moudgalya
  and Motrunich~\cite{moudgalya2022}, the invariant subspaces are the blocks of
  the algebra generated by the local terms (or gates). Its commutant has
  exponential dimension, and the product-basis Krylov sectors (``classical''
  fragmentation) can be refined further by quantum
  fragmentation~\cite{moudgalya2022}. Level statistics are meaningful only
  inside one such block. Pooling blocks measures their number, not their
  chaoticity.
\end{itemize}

\subsection{Eigenstate entanglement}

For a random pure state on $\mathbb{C}^m\otimes\mathbb{C}^n$ with $m\le n$,
Page's mean entanglement entropy is $S_\mathrm{Page}\approx\ln m-m/(2n)$, short
of the maximum by about half a unit of information~\cite{page1993}. Eigenstates
of chaotic Hamiltonians approach this value in the middle of the spectrum.
Conservation laws lower it: Vidmar and Rigol showed that for random states at
fixed particle number away from half filling, the average entropy falls short of
the maximum by an amount growing as the square root of the
volume~\cite{vidmar2017}. For a fragmented system the right benchmark is
therefore a random state \emph{inside the relevant block}, not $S_\mathrm{Page}$
of the full space.

Volume-law eigenstates signal thermalisation. Area-law or otherwise anomalously
weakly entangled eigenstates embedded in a thermal spectrum are the defining
signature of quantum many-body scars (QMBS). QMBS were identified as the origin
of the persistent revivals in the Rydberg-atom experiment~\cite{bernien2017},
via scarred eigenstates of the PXP model~\cite{turner2018}. Embedding
constructions produce them systematically~\cite{shiraishi2017}. Reviews cover
the semiclassical picture~\cite{serbyn2021}, exact
results~\cite{moudgalya2022review}, and the quasiparticle
perspective~\cite{chandran2023}. Scars also exist intrinsically in driven
systems: Sugiura, Kuwahara and Saito proved that a PXP-type Floquet model hosts
exact scar eigenstates of $U$ at specific driving periods, a rigorous breakdown
of Floquet ETH~\cite{sugiura2021}.

\emph{Degeneracies make single-eigenvector entanglement ill-defined.} If an
eigenvalue of $U$ is $g$-fold degenerate, any orthonormal basis of the
eigenspace is equally valid, and the entanglement of ``the'' eigenvectors
depends on that arbitrary choice. Earlier R8 found that 43--74\,\% of the stored
rule-150 eigenstates lie in such eigenspaces, with multiplicity up to 35. This
report therefore (i)~resolves all symmetries first, so the remaining basis
freedom is as small as possible, (ii)~records $g$ for every state, and (iii)~adds
a basis-independent quantity for every eigenspace: the Haar-averaged R\'enyi-2
entropy (Sec.~\ref{sec:methods}).

\subsection{Quantum circuits and automata in particular}

Brickwork Floquet circuits are the setting in which many-body quantum chaos has
been connected analytically to random-matrix theory. Kos, Ljubotina and Prosen
derived RMT spectral correlations for periodically kicked many-body
systems~\cite{kos2018}. Chan, De Luca and Chalker solved a Floquet circuit of
random two-site gates (depth two, like ours) in the large-$q$ limit, including
its spectral form factor and entanglement growth~\cite{chan2018}. The circuits
studied here are the opposite extreme: \emph{one fixed} controlled gate,
repeated.

Automaton-based circuits are a known source of non-ergodic eigenstates.
Gopalakrishnan and Zakirov constructed facilitated quantum cellular automata,
mostly from Clifford gates. Their eigenstate entanglement differs from generic
chaotic systems in all five models, and some models have exponentially many
eigenstates with disentangled sites, tied to reflection
symmetry~\cite{gopalakrishnan2018}. Iadecola and Vijay showed that deforming a
classical reversible cellular automaton gives generically non-integrable Floquet
systems that inherit exact non-thermal eigenstates from the automaton's periodic
orbits~\cite{iadecola2020}. The automata themselves (our $V=X$ circuits) include
the interacting integrable Rule~54~\cite{buca2021} and the Floquet-PXP Rule~201,
whose solitons are domain walls between three vacua~\cite{wilkinson2020}.

\subsection{Two structural facts used throughout}

\paragraph{$U$ is real orthogonal and a product of two reflections.} Both gates,
$H$ and $X$, are real symmetric involutions, so every controlled gate is one and
each layer $L_{A,B}$ is a real symmetric involution (a reflection). Hence
$U=L_BL_A$ is real orthogonal. Two consequences:
\begin{enumerate}
\item Complex conjugation $K$ commutes with $U$, so within every
  symmetry-resolved block the spectrum is symmetric under
  $\theta\to-\theta$. The independent spectral data are the $|\theta|\in[0,\pi]$.
  Ratios computed over the full circle would count each spacing twice and add
  spurious small gaps where mirror pairs approach $0$ and $\pi$.
\item By the two-subspace theorem~\cite{jordan1875,halmos1969}, the space
  decomposes into the intersections of the $\pm1$ eigenspaces of $L_A$ and
  $L_B$ (eigenvalues $\pm1$ of $U$) and two-dimensional subspaces on which $U$
  rotates by $\pm2\varphi_j$, where the $\varphi_j$ are the principal angles
  between the $+1$ eigenspaces of the two layers. For generic reflections these
  angles follow a real ($\beta=1$) Jacobi ensemble, whose bulk statistics are
  those of the GOE. We verified the reference value directly with the same
  routine used for the data: products of two Haar-random reflections of size
  $300$--$800$ give $\langle r\rangle=0.533\pm0.004$ on the half spectrum,
  compared with $0.535\pm0.005$ for a GOE control and $0.604\pm0.005$ for a GUE
  control.
\end{enumerate}
In addition $T=L_AK$ is an antiunitary with $TUT^{-1}=U^{-1}$ and $T^2=1$, a
time-reversal symmetry, and also points to the orthogonal class. The chaotic
reference is therefore $\langle r\rangle_\mathrm{GOE}=0.5307$ for every block that
$T$ maps to itself: all blocks with open boundaries, and the self-conjugate momentum
sectors on a ring. In a ring-momentum sector $j\ne-j$, $T$ maps the sector to its
partner $-j$, no antiunitary symmetry remains inside it, and the chaotic reference is
the CUE value $0.5996$.

\paragraph{The X gate makes $U$ a permutation.} With $V=X$ every basis state is
mapped to a single basis state. On a cycle of length $\ell$ the spectrum is the
$\ell$-th roots of unity, a picket fence with $r\equiv1$, and the eigenvectors
are Fourier modes along the cycle. Each Krylov sector of a permutation is a
single cycle, so this holds sector by sector. The commutant study R-T26 found
exactly this, $P(r)=\delta(r-1)$, for rule 156. Neither Poisson nor RMT applies,
and an X-gate $\langle r\rangle$ measures orbit arithmetic, not chaos.
